% ===========================================================================
%  Cuentas Nacionales de Transferencia para Mexico, 2016-2024
%
%  Compilar:  pdflatex paper -> bibtex paper -> pdflatex paper x2
%
%  Las figuras (figures/) y las tablas (tables/) las genera paper.ipynb.
%  No editarlas a mano: se sobrescriben en cada corrida del cuaderno.
% ===========================================================================
\documentclass[11pt,a4paper]{article}
\usepackage[utf8]{inputenc}
\usepackage[T1]{fontenc}
\usepackage[spanish,es-nodecimaldot]{babel}
\usepackage{lmodern}
\usepackage{amsmath,amssymb}
\usepackage{graphicx}
\usepackage{booktabs}
\usepackage{array}
\usepackage{tabularx}
\usepackage{xcolor}
\usepackage{caption}
\usepackage{float}
\usepackage{placeins}
\usepackage[margin=2.6cm]{geometry}
\usepackage[round,authoryear]{natbib}
\usepackage[hidelinks]{hyperref}
\graphicspath{{figures/}}
\captionsetup{font=small,labelfont=bf}
\setlength{\parskip}{0.35em}
\title{\textbf{Cuentas Nacionales de Transferencia para México, 2016--2024}\\[0.3em]
\large Construcción y conciliación de una base persona-año a partir de la ENIGH
y el Sistema de Cuentas Nacionales}
\author{Juan Pablo López Reynosa \and Héctor Villarreal}
\date{\today}
\begin{document}
\maketitle
\begin{abstract}
Este documento construye una base persona-año de Cuentas Nacionales de Transferencia
para México a partir de cinco levantamientos bienales de la Encuesta Nacional de
Ingresos y Gastos de los Hogares (ENIGH) 2016--2024, con 1.46 millones de registros,
conciliada con el Sistema de Cuentas Nacionales mediante un único factor de
reescalamiento por variable y año. El procedimiento toma de la encuesta la forma del
perfil por edad y de las cuentas nacionales el nivel, garantiza el cierre contra siete
agregados de control y reporta la brecha entre los dos métodos independientes de cálculo
del déficit del ciclo de vida en lugar de forzarla a cero.

Los resultados documentan tres asimetrías estructurales. Primera: el déficit del ciclo
de vida en México es un resultado de composición por sexo. El perfil masculino alcanza
superávit entre los 33 y los 62 años en 2024; el femenino no lo alcanza en ninguna edad
de ninguno de los cinco levantamientos. Las mujeres concentran el 76.3\,\% del déficit
agregado. Segunda: el déficit no crece en términos agregados --- bajó de 47.4\,\% a
42.2\,\% del consumo entre 2016 y 2024 ---, pero se desplaza hacia la vejez: el déficit
per cápita de los mayores de 60 años subió 71.7\,\% en términos reales mientras su peso
poblacional pasó de 11.3\,\% a 15.1\,\%. Tercera: los dos extremos del ciclo se
financian por vías opuestas. La niñez cubre su déficit en 88.6\,\% con recursos privados;
la vejez lo cubre en 54.5\,\% con transferencias públicas netas. Un envejecimiento
demográfico que aumente el peso de la vejez se traduce de forma directa en presión
presupuestaria, mientras un cambio en el peso de la niñez se absorbe dentro de los
hogares. El sistema público no compensa la brecha de financiamiento por sexo en la vejez:
las mujeres mayores de 60 años reciben 24.0\,\% menos en transferencias públicas netas
que los hombres, pese a que su déficit duplica al de ellos.
\end{abstract}
% ---------------------------------------------------------------------------
\section{Objetivo}
% ---------------------------------------------------------------------------
Este documento construye una base persona-año para México que permite calcular perfiles
económicos por edad bajo el marco de Cuentas Nacionales de Transferencia (NTA) y los concilia
con el Sistema de Cuentas Nacionales de México. El producto es una base de 1.46 millones de
registros persona-año que cubre cinco levantamientos bienales de la Encuesta Nacional de
Ingresos y Gastos de los Hogares (ENIGH) --- 2016, 2018, 2020, 2022 y 2024 --- y de la cual se
derivan perfiles por año y edad exacta.
El ejercicio resuelve un problema de compatibilidad entre dos fuentes con propósitos distintos.
La encuesta observa la estructura por edad del ingreso y del gasto, pero subcapta el nivel: los
factores de reescalamiento de la sección~\ref{sec:reescalamiento} van de 1.9 a 9.9 según la
variable y el año, es decir, la encuesta capta entre el 10\,\% y el 53\,\% del agregado
correspondiente de cuentas nacionales. Las cuentas nacionales fijan el nivel pero carecen de
dimensión etaria. El procedimiento toma de la encuesta la forma por edad y de las cuentas
nacionales el nivel, mediante un reescalamiento proporcional por variable y año.
El documento persigue tres objetivos operativos. El primero es dejar explícita la identidad
contable del ciclo de vida y calcularla por dos vías independientes, con reporte de la brecha
entre ambas en lugar de un cierre forzado. El segundo es documentar cada imputación ---
contribuciones a la seguridad social, impuesto sobre la renta, alquiler imputado de la vivienda
propia, consumo de gobierno y transferencias dentro del hogar --- con su ecuación, su agregado
de control y su supuesto. El tercero es acotar qué parte del resultado proviene del dato y qué
parte de la conciliación.
% ---------------------------------------------------------------------------
\section{Revisión de literatura}
% ---------------------------------------------------------------------------
\subsection{NTA como metodología}
Las Cuentas Nacionales de Transferencia extienden el sistema de cuentas nacionales e incorporan
la edad como dimensión de la contabilidad macroeconómica \citep{lee_mason_2011,
un_nta_manual_2013}. El objeto de medición es el ciclo de vida económico: en la niñez y en la
vejez el consumo excede al ingreso laboral, y esa diferencia se financia con transferencias
públicas, transferencias privadas y reasignaciones por activos. El marco cuantifica esos flujos
por edad y exige que sus agregados reproduzcan las cuentas nacionales del país, de modo que los
perfiles etarios resulten comparables entre economías y consistentes con el sistema de cuentas
nacionales \citep{sna_2008}.
La construcción sigue un orden fijo. Se define un agregado de control por variable desde las
cuentas nacionales; se estima la forma por edad desde encuestas de hogares y registros
administrativos; y se reescala la forma micro para que su agregado reproduzca el control macro.
El manual de Naciones Unidas \citep{un_nta_manual_2013} formaliza el procedimiento y fija las
convenciones de clasificación que hacen comparables los resultados entre países. La literatura
posterior ha usado los perfiles resultantes para medir el dividendo demográfico y la
sostenibilidad de los sistemas de transferencias públicas \citep{mason_lee_2007, lee_1994}.
\subsection{Casos especiales}
Dos imputaciones concentran la mayor parte de la modelización propia de una aplicación NTA y
merecen tratamiento separado. La primera es el alquiler imputado de la vivienda propia. La
encuesta capta el gasto como desembolso: un hogar arrendatario paga renta y ese pago se
registra, mientras que un hogar propietario consume el mismo servicio habitacional sin
desembolso, de modo que el servicio no aparece por ningún lado. El consumo final de los hogares
de las cuentas nacionales sí lo incluye. El marco NTA resuelve la asimetría con el
reconocimiento de que el alquiler imputado es dos cosas a la vez --- por el lado del consumo, el
servicio habitacional que el hogar disfruta; por el lado del ingreso de activos, el rendimiento
del activo vivienda --- y lo incorpora a ambos lados de la identidad, ecuaciones~\eqref{eq:cpriv}
y~\eqref{eq:ykalq}.
La segunda es la asignación a personas del gasto observado a nivel hogar. La encuesta registra
alimentos, vivienda y servicios como gasto del hogar, no de sus miembros, y el perfil etario del
consumo depende por completo de cómo se reparta ese gasto. La alternativa habitual consiste en
imponer una escala de equivalencia externa, construida para otro país o para otro propósito. El
procedimiento que se documenta aquí extrae la escala del propio dato con un algoritmo de punto
fijo, ecuaciones~\eqref{eq:alloc1}--\eqref{eq:allocconv}: parte de un reparto per cápita, estima
el perfil etario implícito, redistribuye el gasto del hogar según ese perfil y repite hasta la
convergencia, con conservación exacta del total del hogar en cada iteración. La alternativa
econométrica sería un estimador del tipo \citet{deaton_paxson_1998}, que identifica economías de
escala del hogar a partir de la respuesta del gasto en alimentos al tamaño del hogar; queda
anotada como contraste externo del perfil de consumo por edad.
\subsection{Aplicaciones y usos de NTA en América Latina}
La evidencia aplicada en América Latina ha usado el marco sobre todo para medir cómo el
envejecimiento altera las transferencias entre generaciones y presiona los sistemas de protección
social. \citet{RoseroBixby2011Generational} compara Brasil, Chile, Costa Rica, México y Uruguay, y
encuentra que el envejecimiento eleva la carga de sostén sobre las edades activas y vuelve
visibles las tensiones entre el circuito familiar y el público. \citet{Mason2011Population}
sintetizan resultados comparados y subrayan que el aumento de la población mayor reduce la razón
de soporte y eleva la relevancia fiscal de pensiones, salud y cuidados de largo plazo. Este
documento se ubica en esa tradición, pero su unidad de análisis no es la proyección sino la
construcción: qué parte de un perfil por edad proviene de la encuesta y qué parte de la
conciliación con las cuentas nacionales.
Una segunda línea muestra que los promedios nacionales ocultan diferencias distributivas grandes,
y es la que más se acerca a los resultados de la sección~\ref{sec:resultados}. Para Ecuador,
\citet{Rosero2024Socioeconomic} encuentra que algunas transferencias públicas son progresivas
mientras las pensiones de retiro resultan fuertemente regresivas. Para Perú,
\citet{Olivera2023distributive} documenta que los cambios distributivos del periodo favorecieron
relativamente más a los grupos jóvenes y con mayor escolaridad. La implicación metodológica es
directa: un perfil promedio por edad puede ser consistente con el agregado macro y aun así
describir mal a los subgrupos que lo componen.
Para México, \citet{MejiaGuevaraRivero2024} muestran que el análisis del ingreso laboral y del
consumo por edad cambia de forma sustantiva al desagregar por grupo socioeconómico y por sexo, y
al incorporar el trabajo no remunerado mediante cuentas de transferencias de tiempo. Sus
resultados indican que las diferencias no se explican sólo por edad, sino por la interacción entre
escolaridad, género y organización del trabajo pagado y no pagado. Este documento comparte ese
diagnóstico y aporta evidencia complementaria por el lado del financiamiento: la
sección~\ref{sec:resultados} muestra que el ciclo de vida masculino alcanza superávit en un tramo
amplio de edades mientras el femenino no lo alcanza en ninguna, de modo que el agregado nacional
que aparece deficitario en todo el ciclo es un resultado de composición.
\subsection{Identidad contable completa del ciclo de vida}
\label{sec:identidad}
La restricción presupuestaria de una persona de edad $a$ en el año $t$ iguala los recursos con
los usos:
\begin{equation}
\label{eq:presupuesto}
\underbrace{Y_L(a,t) + Y_K(a,t)}_{\text{ingreso propio}} \;+\;
\underbrace{\tau^{+}(a,t)}_{\text{transf. recibidas}}
\;=\;
\underbrace{C(a,t)}_{\text{consumo}} \;+\; \underbrace{S(a,t)}_{\text{ahorro}} \;+\;
\underbrace{\tau^{-}(a,t)}_{\text{transf. pagadas}}
\end{equation}
Al reordenar los términos, el exceso de consumo sobre ingreso laboral queda igualado a la suma
de las transferencias netas y de la reasignación por activos:
\begin{equation}
\label{eq:reordenada}
\underbrace{C(a,t) - Y_L(a,t)}_{\text{Déficit del Ciclo de Vida}}
\;=\;
\underbrace{\big(\tau^{+}(a,t) - \tau^{-}(a,t)\big)}_{\text{transferencias netas}}
\;+\;
\underbrace{\big(Y_K(a,t) - S(a,t)\big)}_{\text{reasignación por activos}}
\end{equation}
\begin{equation}
\label{eq:lcd}
LCD(a,t) \;\equiv\; C(a,t) - Y_L(a,t)
\end{equation}
La convención de signo importa para la lectura de cualquier resultado. El déficit del ciclo de
vida se define como $C - Y_L$ y no como $Y_L - C$: en la niñez y en la vejez el consumo excede al
ingreso laboral y el déficit es positivo, mientras que en la edad activa el ingreso laboral
excede al consumo y el saldo es negativo, es decir, superávit.
Los dos lados de la ecuación~\eqref{eq:reordenada} dan lugar a dos formas independientes de
calcular la misma cantidad, y este documento implementa ambas en paralelo:
\begin{equation}
\label{eq:m1}
LCD_{M1}(a,t) = C(a,t) - Y_L(a,t)
\end{equation}
\begin{equation}
\label{eq:m2}
LCD_{M2}(a,t) = \big(\tau^{+}(a,t) - \tau^{-}(a,t)\big) + \big(Y_K(a,t) - S(a,t)\big)
\end{equation}
El primer método mide el déficit; el segundo mide cómo se financia. En una contabilidad cerrada
ambos coinciden. Aquí no coinciden de forma exacta, y la diferencia se reporta en lugar de
forzarse a cero: es consecuencia directa de que las cuentas nacionales no publiquen un agregado
de transferencias privadas netas, y su magnitud es derivable de forma analítica
(sección~\ref{sec:validacion}). El segundo método se reporta en su versión completa, que incluye
el bloque privado observado:
\begin{equation}
\label{eq:m2full}
LCD^{\,c}_{M2}(a,t) = \big(\tau^{+}_{pub}(a,t) - \tau^{-}_{pub}(a,t)\big)
+ \tau^{+}_{ext}(a,t) + \tau^{priv}(a,t) + \big(Y_K(a,t) - S(a,t)\big)
\end{equation}
El último término desempeña el papel de $(Y_K - S)$ de la ecuación~\eqref{eq:m2}, porque su
agregado de control es la reasignación neta por activos y no el ingreso de activos bruto. No
existe, por tanto, un término de ahorro separado.
% ---------------------------------------------------------------------------
\section{Manejo de datos}
% ---------------------------------------------------------------------------
\subsection{Fuentes macro}
Los agregados de control provienen del Sistema de Cuentas Nacionales de México, cuentas por
sectores institucionales \citep{inegi_scnm}, en pesos corrientes de cada año, y se deflactan a precios de 2024. Las series
utilizadas son la remuneración de asalariados, el excedente bruto de operación por cuenta propia
y el excedente bruto de explotación; el consumo de capital fijo, el saldo neto de rentas de la
propiedad y el ahorro neto; el gasto de consumo final de los hogares y del gobierno, con su
apertura de educación y salud; las prestaciones sociales, las transferencias sociales en especie
y las transferencias netas del exterior; y los impuestos corrientes sobre el ingreso junto con
las contribuciones sociales efectivas de los empleadores. El Anexo~\ref{anexo:series} identifica
cada serie con su clave del sistema de cuentas y el agregado de control que alimenta.
El agregado de transferencias públicas pagadas suma los impuestos corrientes al ingreso y las
contribuciones sociales de los empleadores, efectivas e imputadas. Las imputadas son la
contraparte de las prestaciones que un empleador paga de forma directa, sin fondo constituido, y
representan entre 5\,\% y 8\,\% de ese agregado según el año. Se incluyen por simetría contable:
el agregado del ingreso laboral las contiene dentro de la remuneración de asalariados, y la
prestación que financian ya forma parte de las prestaciones sociales recibidas, de modo que
excluirlas sólo del lado del pago haría aparecer al hogar cobrando la contribución y su prestación
sin pagarla nunca.
Sobre las series publicadas se derivan tres agregados que no aparecen como tales. El ingreso de
activos ampliado incorpora el saldo neto de rentas de la propiedad frente al resto del mundo, y
de él se descuenta el ahorro neto:
\begin{equation}
\label{eq:activos}
Y_A(t) = Y_K(t) + D^{prop}(t), \qquad
\big(Y_K - S\big)(t) \;\equiv\; Y_A(t) - S(t)
\end{equation}
\begin{equation}
\label{eq:cfuncion}
C_{edu}(t) = C^{hog}_{edu}(t) + C^{gob}_{edu}(t), \qquad
C_{sal}(t) = C^{hog}_{sal}(t) + C^{gob}_{sal}(t)
\end{equation}
\subsection{Deflactación}
Todos los montos monetarios de este documento están expresados en pesos constantes de 2024. Sea
$\pi(t)$ el deflactor implícito con $\pi(2024)=1$; cada monto micro y cada agregado de control se
divide entre el deflactor de su año:
\begin{equation}
\label{eq:deflactor}
x^{r}_i(t) = \frac{x_i(t)}{\pi(t)}, \qquad
V^{r}_{macro}(t) = \frac{V_{macro}(t)}{\pi(t)}
\end{equation}
El deflactor acumula 48.3\,\% entre 2016 y 2024, con incrementos bienales de 12.1\,\%, 9.1\,\%,
11.5\,\% y 8.7\,\%. La deflactación se aplica antes del reescalamiento de la
sección~\ref{sec:reescalamiento} y no lo altera: el factor de reescalamiento es un cociente entre
un agregado macro y un agregado micro del mismo año, de modo que ambos se dividen entre el mismo
$\pi(t)$ y el cociente queda idéntico. La deflactación cambia el nivel de las cifras reportadas,
no la conciliación.
\subsection{Fuente micro}
La fuente micro es la Encuesta Nacional de Ingresos y Gastos de los Hogares del INEGI
\citep{inegi_enigh}, levantamientos 2016, 2018, 2020, 2022 y 2024. Las tablas de cada
levantamiento --- población, ingresos, gasto por persona, gasto del hogar, concentrado del hogar
y trabajos --- se consolidan en una fila por persona y año, con montos en pesos
trimestrales. El Cuadro~\ref{tab:muestra} resume la composición de cada levantamiento. La muestra
va de 257\,805 a 315\,743 registros y la población expandida de 120.9 a 130.3 millones de
personas. El envejecimiento del periodo es visible en la propia muestra: la edad media sube de
31.0 a 34.4 años, la participación de los menores de 15 años baja de 27.2\,\% a 22.3\,\% y la de
los mayores de 60 sube de 11.3\,\% a 15.1\,\%.
\input{tables/tab_muestra.tex}
La encuesta capta un trimestre de referencia y las cuentas nacionales son anuales, de modo que
toda comparación entre ambas exige una regla de cierre explícita. Cualquier total anual de este
documento se construye como
\begin{equation}
\label{eq:cierre}
X(t) \;=\; \Big(\sum_{i \in t} x_i \cdot f_i\Big) \cdot 4
\end{equation}
donde $f_i$ es el factor de expansión de la persona $i$. El factor de expansión no se modifica en
ninguna etapa del procedimiento: la conciliación actúa sobre montos monetarios, no sobre pesos
muestrales.
Esa regla implica un supuesto: el trimestre de referencia se toma como representativo del año.
El nivel anual no depende de él, porque queda fijado por el agregado
de control, pero la forma por edad sí lo haría si existiera estacionalidad correlacionada con la
edad --- gasto escolar concentrado al inicio del ciclo lectivo, ingreso agrícola concentrado en la
cosecha. El levantamiento ocurre siempre en el mismo periodo del año, de modo que el sesgo, de
existir, es estable entre levantamientos y no contamina la comparación temporal.
\subsection{Correspondencia entre variables micro y agregados de control}
La correspondencia entre cada variable de la base persona-año y su agregado de control está en
el Anexo~\ref{anexo:mapeo}. Es la única declaración del mapa de escalamiento: una variable que no
aparece en el Cuadro~\ref{tab:scale_map} no recibe factor de reescalamiento. El
Cuadro~\ref{tab:excluidas} documenta las exclusiones deliberadas y la razón de cada una, que no
es la misma en todos los casos.
% ---------------------------------------------------------------------------
\section{Definición de variables}
\label{sec:variables}
% ---------------------------------------------------------------------------
\subsection{Ingreso laboral}
El ingreso laboral individual suma el ingreso subordinado y la componente laboral del ingreso
mixto. La encuesta reporta las percepciones por clave, y la clasificación separa el ingreso
subordinado, el ingreso mixto del trabajador por cuenta propia y el ingreso de capital puro. El
ingreso mixto retribuye a la vez al trabajo y al capital de quien trabaja por cuenta propia, y la
convención del manual de Naciones Unidas lo divide en dos tercios laborales y un tercio de capital
\citep{un_nta_manual_2013}; la división 2/3--1/3 se registra como supuesto E1 en el
anexo~\ref{anexo:supuestos}. La misma convención se aplica del lado macro, de modo que la
proporción es idéntica en el numerador y en el denominador del factor de reescalamiento: el
agregado de control suma la remuneración de asalariados --- con las contribuciones sociales de los
empleadores, efectivas e imputadas --- y dos tercios del ingreso mixto neto de su depreciación,
ecuación~\eqref{eq:ylmacro}.
\begin{equation}
\label{eq:yl1}
Y^{sub}_i = \sum_{k \in K_{sub}} y_{ik}, \qquad
Y^{mix}_i = \sum_{k \in K_{mix}} y_{ik}, \qquad
Y^{cap}_i = \sum_{k \in K_{cap}} y_{ik}
\end{equation}
\begin{equation}
\label{eq:yl2}
Y^{mix,L}_i = \tfrac{2}{3}\, Y^{mix}_i, \qquad Y^{mix,K}_i = \tfrac{1}{3}\, Y^{mix}_i
\end{equation}
\begin{equation}
\label{eq:yl3}
Y_{L,i} = Y^{sub}_i + Y^{mix,L}_i, \qquad
Y_{K,i} = Y^{cap}_i + Y^{mix,K}_i
\end{equation}
\begin{equation}
\label{eq:ylmacro}
Y_L(t) = \underbrace{W(t)}_{\text{remuneración de asalariados}}
+ \tfrac{2}{3}\,\big[M(t) - \delta_M(t)\big]
\end{equation}
\subsection{Consumo}
El consumo se mide a nivel persona, lo que obliga a resolver que la encuesta observe buena parte
del gasto a nivel hogar. La asignación usa un algoritmo iterativo de punto fijo que construye su
propia escala de equivalencia etaria a partir del dato en vez de importar una escala externa:
parte de un reparto per cápita, ecuación~\eqref{eq:alloc1}; estima la participación de cada edad
en el agregado, ecuación~\eqref{eq:alloc2}; redistribuye el gasto del hogar en proporción a esas
participaciones, ecuación~\eqref{eq:alloc3}; y repite hasta que se cumple el criterio de
convergencia, ecuación~\eqref{eq:allocconv}. El reparto conserva el total del hogar de forma
exacta en cada iteración, ecuación~\eqref{eq:alloccons}, de modo que ninguna configuración del
algoritmo puede mover un agregado. El consumo de gobierno se imputa de forma plana per cápita,
ecuación~\eqref{eq:cgob}, y se incorpora \emph{después} de las transferencias dentro del hogar:
el consumo público se financia con impuestos y no con transferencias privadas, de manera que su
incorporación previa lo convertiría en un déficit que algún miembro del hogar tendría que cubrir.
El reescalamiento de la sección~\ref{sec:reescalamiento} trata este bloque de forma partida: el
consumo de gobierno conserva el valor con que se calibró y el factor se aplica únicamente al
residuo privado, cuyo agregado de control es el consumo total menos el consumo de gobierno,
ecuación~\eqref{eq:cparticion}. Sin esa partición el factor volvería a escalar una componente ya
calibrada y la dejaría al doble de su nivel de control.
La imputación es plana para el total del consumo de gobierno y no separa su componente colectivo
--- defensa, administración, orden público --- del individualizable, que es cerca de la mitad y
corresponde a educación y salud públicas. Darle a esa mitad la forma etaria del gasto privado en
educación y salud parecería una mejora, pero descansa en una equivalencia que la evidencia no
sostiene: en salud, el gasto de bolsillo puede ser \emph{inverso} al beneficio público recibido,
porque quien desembolsa lo hace a menudo por no haber resuelto su necesidad en el sistema público,
mientras que quien registra gasto nulo pudo haberse atendido en él. Sin una fuente de utilización
efectiva de servicios públicos por edad, imponer el perfil del gasto privado arriesga invertir el
signo de la asignación. La imputación plana no inventa una forma que no se ha medido, y es la
opción conservadora mientras no se incorporen registros de matrícula y de utilización.
\begin{equation}
\label{eq:alloc1}
x_i^{(1)} = \frac{X_h}{n_h}
\end{equation}
\begin{equation}
\label{eq:alloc2}
T_a^{(k)} = \sum_{i : a_i = a} x_i^{(k)} \cdot f_i,
\qquad
\theta_a^{(k)} = \frac{T_a^{(k)}}{\sum_{a'} T_{a'}^{(k)}}
\end{equation}
\begin{equation}
\label{eq:alloc3}
x_i^{(k+1)} = X_h \cdot \frac{\theta_{a_i}^{(k)}}{\sum_{j \in h} \theta_{a_j}^{(k)}}
\end{equation}
\begin{equation}
\label{eq:alloccons}
\sum_{i \in h} x_i^{(k)} = X_h \qquad \text{para toda iteración } k
\end{equation}
\begin{equation}
\label{eq:allocconv}
\sum_a \big|\theta_a^{(k)} - \theta_a^{(k-1)}\big| < \varepsilon
\end{equation}
\begin{equation}
\label{eq:cgob}
C^{gob}_i(t) = \frac{C^{gob}(t)}{\big(\sum_i f_i(t)\big)\cdot 4}
\end{equation}
\begin{equation}
\label{eq:cparticion}
C^{adj}_i(t) = \big(C_i(t) - C^{gob}_i(t)\big)\cdot\lambda^{priv}_C(t) + C^{gob}_i(t),
\qquad
\lambda^{priv}_C(t) = \frac{C_{total}(t) - C^{gob}(t)}
{\sum_i \big(C_i(t) - C^{gob}_i(t)\big)\cdot f_i \cdot 4}
\end{equation}
\subsection{Transferencias fiscales pagadas}
La encuesta reporta ingreso neto, no carga fiscal, de modo que el vector de transferencias
pagadas se reconstruye persona por persona a partir de la normativa. Las contribuciones a la
seguridad social se calculan sobre el salario base de cotización, con piso de una unidad de
medida, topes diferenciados por régimen, y valor nulo para los informales; el pago que entra a
la contabilidad es la suma de la contribución patronal y la de gobierno, ecuación~\eqref{eq:css},
y excluye de forma deliberada la contribución del trabajador, que ya viene descontada del ingreso
observado y cuya suma sería doble conteo. El impuesto sobre la renta exige resolver el problema
inverso: la tarifa se aplica sobre el ingreso bruto mientras la encuesta observa el neto, y la
solución es una iteración de punto fijo, ecuaciones~\eqref{eq:isr1}--\eqref{eq:isrconv}, cuya
convergencia está garantizada porque la tasa marginal máxima es 0.35 y el operador es por tanto
una contracción. Esa rutina cubre únicamente el componente asalariado, mientras que el agregado
de control abarca todas las rentas, incluidas las de personas morales; para cerrar esa asimetría
se añade el tramo de personas morales sobre el ingreso de capital observado con la tasa del
artículo 9 de la Ley del ISR, ecuación~\eqref{eq:isrpm}, y ambos bloques se consolidan en la
ecuación~\eqref{eq:taupagada}. La base gravable de ese tramo excluye el alquiler imputado de la
vivienda propia, que es un ingreso no monetario y no constituye hecho gravable; es el mismo
criterio del simulador fiscal del CIEP, que construye la base de personas morales como el ingreso
de capital menos la imputación de vivienda. Como la tasa es plana, el resultado no depende de la
frecuencia de la base.
Toda la imputación fiscal opera sobre base anual. La encuesta capta un trimestre, mientras que el
salario base de cotización, su piso de una unidad de medida, sus topes por régimen y la tarifa
progresiva son parámetros anuales; aplicarlos al flujo trimestral comprime a casi toda la
población formal contra el piso y contra los primeros tramos de la tarifa. La base se anualiza
antes de aplicar cualquier parámetro y el resultado regresa a periodicidad trimestral, que es la
unidad del resto del ejercicio. Con esta convención la proporción de trabajadores formales con
impuesto positivo va de 65.8\,\% a 92.4\,\% según el año, y la participación del decil superior
en el impuesto imputado queda entre 52.1\,\% y 65.6\,\%.
\begin{equation}
\label{eq:css}
\tau^{ss}_i = \tau^{ss,pat}_i + \tau^{ss,gob}_i
\end{equation}
\begin{equation}
\label{eq:isr1}
g(B) = q_k + (B - l_k)\cdot t_k, \qquad \text{para } B \in (l_k, u_k]
\end{equation}
\begin{equation}
\label{eq:isr2}
T(y) = y^{neto}_i + \max\Big( g\big(B_i(y)\big) - se\big(B_i(y)\big),\; 0 \Big),
\qquad B_i(y) = \max\big(y - \tau^{ss,tra}_i,\,0\big)
\end{equation}
\begin{equation}
\label{eq:isrconv}
\max_i \big| y^{(m+1)}_i - y^{(m)}_i \big| < 1 \text{ peso}
\end{equation}
\begin{equation}
\label{eq:isrpm}
\tau^{isr}_i = \tau^{isr,asal}_i + t_{PM}\cdot \max\big(Y^{obs}_{K,i},\,0\big),
\qquad t_{PM} = 0.30
\end{equation}
\begin{equation}
\label{eq:taupagada}
\tau^{-}_{pub,i} = \tau^{ss}_i + \tau^{isr}_i
\end{equation}
\subsection{Transferencias recibidas}
Las transferencias públicas recibidas se anclan a la suma de las prestaciones sociales y las
transferencias sociales en especie, ecuación~\eqref{eq:taurecibida}. Las transferencias del
exterior se anclan a las transferencias netas del exterior y no al bloque de cooperación
internacional corriente, y la razón es de orden de magnitud: las remesas efectivas van de 0.5 a
1.18 billones de pesos según el año, mientras que ese segundo bloque ronda los 20 mil millones,
es decir, entre 25 y 60 veces menos. Como la serie de transferencias netas ya viene neta de los
pagos de los hogares al exterior, el flujo pagado se fija en cero y el
neto coincide con el recibido.
Este bloque tiene el factor de reescalamiento más alto del ejercicio y merece una advertencia
específica. La encuesta capta entre 10.1\,\% y 12.0\,\% del agregado de control, sobre una base
de 1.6\,\% de las personas y 4.4\,\% de los hogares. La subcaptación de remesas ocurre en dos
márgenes --- cuántos hogares reciben y cuánto recibe cada uno --- y el reescalamiento proporcional
atribuye la corrección entera al segundo: el monto medio por receptor pasa de 16\,088 a
145\,339 pesos trimestrales, es decir 581\,354 pesos anuales, una cifra que excede lo que
documentan las estadísticas de envíos. La sección~\ref{sec:validacion} reporta los resultados que
dependen de este canal con y sin el ajuste.
\begin{equation}
\label{eq:taurecibida}
\tau^{+}_{pub}(t) = B^{soc}(t) + B^{esp}(t)
\end{equation}
\subsection{Asignaciones dentro y entre hogares}
Dentro del hogar existe una reasignación privada que este documento hace visible en lugar de
dejarla implícita en un residuo contable. Para cada persona se calcula el ingreso propio total,
ecuación~\eqref{eq:ytot}, y sus posiciones de déficit y superávit, ecuación~\eqref{eq:defsup};
dentro de cada hogar se emparejan miembros deficitarios con superavitarios mediante una matriz de
pagos sujeta a las restricciones de la ecuación~\eqref{eq:restr}, y el saldo neto por persona,
ecuación~\eqref{eq:taunet}, suma exactamente cero a nivel de hogar, ecuación~\eqref{eq:sumacero}.
A diferencia de un residuo contable, el método puede dejar déficit sin cubrir cuando el hogar no
dispone de superávit suficiente, ecuación~\eqref{eq:defnc}; ese residuo se financia por fuerza
desde fuera del hogar, vía transferencias públicas, remesas o transferencias entre hogares. El
bloque privado, que suma la reasignación dentro del hogar y las transferencias observadas entre
hogares distintos, queda observado y sin reescalar, decisión que genera una brecha conocida y
explicable en el segundo método de cálculo del déficit.
\begin{equation}
\label{eq:ytot}
Y^{tot}_i = Y_{L,i} + Y_{K,i} + \tau^{+}_{pub,i} + \tau^{+}_{ext,i}
\end{equation}
\begin{equation}
\label{eq:defsup}
d_i = \max\big(C_i - Y^{tot}_i,\; 0\big),
\qquad
s_i = \max\big(Y^{tot}_i - C_i,\; 0\big)
\end{equation}
\begin{equation}
\label{eq:restr}
\sum_{k \in h} p_{kj} \le d_j \quad \forall j,
\qquad
\sum_{j \in h} p_{kj} \le s_k \quad \forall k
\end{equation}
\begin{equation}
\label{eq:taunet}
\tau^{in}_i = \underbrace{\sum_{k} p_{ki}}_{\text{recibe}}
- \underbrace{\sum_{j} p_{ij}}_{\text{paga}}
\end{equation}
\begin{equation}
\label{eq:sumacero}
\sum_{i \in h} \tau^{in}_i = 0 \quad \text{para todo hogar } h
\end{equation}
\begin{equation}
\label{eq:defnc}
d^{nc}_i = d_i - \sum_k p_{ki} \; \ge \; 0
\end{equation}
\subsection{Reasignación por activos}
El agregado de control de esta variable no es el ingreso de activos bruto sino la reasignación
neta de ahorro, ecuación~\eqref{eq:activos}. La razón es la ecuación~\eqref{eq:m2}: el segundo
método requiere exactamente el término $(Y_K - S)$, y un agregado ya neto de ahorro evita la
asignación del ahorro por edad, tarea para la cual la encuesta no ofrece información confiable.
En la práctica, el ahorro nunca se asigna por edad de forma independiente y cualquier lectura del
perfil de ingreso de activos debe entenderse como lectura de la reasignación neta. El alquiler
imputado de la vivienda propia entra por este lado además de por el lado del consumo,
ecuación~\eqref{eq:ykalq}, neto de predial y conservación, ecuación~\eqref{eq:alqneto}. El
reparto dentro del hogar es per cápita por los dos lados: cada integrante recibe la misma fracción
del servicio habitacional y del rendimiento del activo. Es el criterio del simulador fiscal del
CIEP, que divide la imputación entre el número de miembros del hogar. La alternativa de asignar
el rendimiento a la persona de referencia concentra el ingreso de activos en las edades donde se
ubica la jefatura del hogar y lo elimina de la niñez, lo que altera la lectura del financiamiento
del déficit infantil sin cambiar ningún agregado.
El alquiler imputado sí tiene contraparte en las cuentas nacionales y está dentro del agregado de
control: el excedente de operación de los hogares, que es la renta imputada de la vivienda propia,
forma parte del excedente bruto de explotación del que se deriva la reasignación por activos. Lo
que no tiene es un agregado de control \emph{propio}, de modo que el factor de reescalamiento se
aplica al bloque completo. Eso distorsiona su composición interna: el alquiler imputado representa
entre 58\,\% y 63\,\% del ingreso de activos observado en la encuesta, frente a una contraparte
macro que ronda entre 10\,\% y 16\,\% del agregado de control, y tras el reescalamiento el bloque
de alquiler queda en torno a 2.8 veces el excedente bruto de operación de los hogares.
La partición que se aplica al consumo no se replica aquí, y la razón es que operaría en sentido
contrario. Antes de cualquier ajuste, el alquiler imputado de la encuesta \emph{excede} a su
contraparte macro --- 1.36 contra 0.90 billones de pesos en 2024 --- porque se construye con un
modelo y no depende de que el informante lo declare, mientras que el ingreso de capital observado
se capta con mucha menor calidad. Anclarlo por separado lo reduciría en un tercio y trasladaría
toda la brecha de captación al capital observado, con un factor cercano a diez sobre una base de
7\,\% de la población, que es la misma configuración que el documento señala como problemática en
el caso de las remesas. A esto se añade un desajuste de concepto: la imputación micro es bruta de
depreciación --- sólo descuenta predial y conservación --- mientras la contraparte macro neta ya la
descuenta, de modo que forzar la igualdad mezclaría un problema de captación con uno de definición.
El bloque se deja sin particionar y la distorsión se reporta como limitación.
\begin{equation}
\label{eq:cpriv}
C^{priv}_i = C^{obs}_i + R^{C}_i
\end{equation}
\begin{equation}
\label{eq:ykalq}
Y_{K,i} = Y^{obs}_{K,i} + R^{K}_i
\end{equation}
\begin{equation}
\label{eq:alqneto}
R^{K}_h = \max\big(R_h - P_h,\; 0\big)
\end{equation}
% ---------------------------------------------------------------------------
\section{Reescalamiento micro a macro}
\label{sec:reescalamiento}
% ---------------------------------------------------------------------------
La conciliación se resuelve con un único escalar por variable y año. Para cada variable $v$ del
mapa de escalamiento se calcula
\begin{equation}
\label{eq:lambda}
\lambda_v(t) \;=\; \frac{V_{macro}(t)}{\displaystyle\sum_i x_{iv}(t)\cdot f_i \cdot 4}
\end{equation}
y se aplica de manera uniforme a todas las observaciones del año:
\begin{equation}
\label{eq:lambdaap}
x^{adj}_{iv}(t) = x_{iv}(t)\cdot \lambda_v(t)
\end{equation}
Tres propiedades justifican esta elección frente a alternativas más elaboradas, como la
calibración por ajuste iterativo proporcional o la reponderación de factores. La primera es el
cierre exacto por construcción: al sustituir~\eqref{eq:lambda} en~\eqref{eq:lambdaap}, el
agregado ajustado reproduce el control macro sea cual sea $x_{iv}$, ecuación~\eqref{eq:cierrelam}.
Esto permite afirmar sin necesidad de un experimento que ninguna imputación previa --- alquiler
imputado, asignación del gasto del hogar, imputación fiscal --- puede romper el cierre agregado.
La segunda es la invariancia de la forma relativa, ecuación~\eqref{eq:invar}: el reescalamiento
preserva toda la información distributiva y etaria de la encuesta, y ningún perfil por edad
cambia de forma, sólo de escala vertical. La tercera es que el factor actúa sobre montos
monetarios y nunca sobre el factor de expansión, de modo que la conciliación no altera la
estructura demográfica de la muestra expandida.
\begin{equation}
\label{eq:cierrelam}
\sum_i x^{adj}_{iv}(t)\cdot f_i\cdot 4 \;=\;
\lambda_v(t)\sum_i x_{iv}(t)\cdot f_i\cdot 4 \;=\; V_{macro}(t)
\end{equation}
\begin{equation}
\label{eq:invar}
\frac{x^{adj}_{iv}}{x^{adj}_{jv}} = \frac{x_{iv}}{x_{jv}}
\end{equation}
El consumo es la única variable que no recibe el factor sobre su total, por la partición ya
descrita en la sección~\ref{sec:variables} (ecuación~\eqref{eq:cparticion}): el bloque de
gobierno conserva su valor calibrado y el factor actúa sólo sobre el residuo privado. El total
sigue cerrando contra el agregado de control del consumo.
El Cuadro~\ref{tab:lambda} del Anexo~\ref{anexo:lambda} reporta los 35 factores, siete variables
por cinco años, con marca de auditoría para los que caen fuera de la banda $[0.5,\,2.0]$. Treinta
y tres de los treinta y cinco quedan fuera de esa banda, y todos por arriba: la encuesta capta
menos que las cuentas nacionales en las siete variables. El rango va de 1.90 en el ingreso
laboral de 2024 a 9.88 en las transferencias del exterior de 2020. Un factor lejos de la unidad
no es en sí mismo un error, sino la medida de la subcaptación de la encuesta; lo que sí
contaminaría la comparación temporal es un factor con cambios grandes entre años sin razón
metodológica, porque en ese caso parte de la variación observada en los perfiles sería artefacto
de la conciliación y no cambio económico real. Los factores son invariantes a la deflactación,
por la razón expuesta en la ecuación~\eqref{eq:deflactor}, de modo que los valores del cuadro son
los mismos en términos nominales y reales.
El costo de este procedimiento define también su alcance. El reescalamiento proporcional atribuye
la totalidad de la discrepancia micro-macro a un error de nivel uniforme, de modo que no corrige
subcaptación diferencial por edad ni por posición en la distribución del ingreso. Donde eso pesa
más es en el ingreso de activos, subcaptado de forma característica en la cola alta, cuyo factor
va de 4.28 a 6.11. El método sirve como control macroeconómico de los perfiles por edad, no como
instrumento para estudiar la desigualdad del ingreso de capital. Cualquier lectura distributiva
de ese bloque excede lo que el procedimiento puede sostener, y el documento no la propone.
% ---------------------------------------------------------------------------
\section{Presentación de resultados}
\label{sec:resultados}
% ---------------------------------------------------------------------------
\subsection{Ingreso laboral, total y por sexo}
El ingreso laboral agregado creció 22.1\,\% en términos reales entre 2016 y 2024, de 12.2 a 14.9
billones de pesos de 2024. El perfil por edad conserva la forma de campana característica del
ciclo laboral: el ingreso es prácticamente nulo antes de los 15 años, se eleva con rapidez entre
los 20 y los 35, y desciende de forma sostenida a partir de los 55. El desplazamiento vertical de
la curva ocurre sin cambio apreciable en el tramo de edades donde el perfil alcanza su máximo, de
modo que el crecimiento del agregado responde al nivel del ingreso y no a un corrimiento del
calendario laboral.
La convergencia por sexo explica buena parte de ese crecimiento. El ingreso laboral real per
cápita de las mujeres subió 25.1\,\% en el periodo, de 15\,016 a 18\,785 pesos trimestrales,
mientras el de los hombres subió 8.6\,\%, de 36\,107 a 39\,222. La razón entre ambos pasó de
41.6\,\% a 47.9\,\%, y de 44.1\,\% a 48.0\,\% si la comparación se restringe al grupo de 25 a 59
años. La brecha se cerró poco más de seis puntos porcentuales sin alterar su orden de magnitud:
al cierre del periodo el ingreso laboral femenino permanece por debajo de la mitad del masculino.
\begin{figure}[H]
\centering
\includegraphics[width=\linewidth]{fig_yl_sexo.pdf}
\caption{Ingreso laboral por edad. Panel izquierdo: total de la población, 2016--2024. Panel
derecho: por sexo de la persona, 2024. Pesos constantes de 2024, trimestrales per cápita.}
\label{fig:yl_sexo}
\end{figure}
\subsection{Consumo y déficit del ciclo de vida}
El consumo agregado creció 11.1\,\% en términos reales entre 2016 y 2024, de 23.2 a 25.8 billones
de pesos de 2024, y superó al ingreso laboral en los cinco levantamientos. El déficit del ciclo de
vida se mantuvo prácticamente estable: pasó de 11.0 a 10.9 billones, una caída real de 1.1\,\%.
Medido como proporción del consumo, el déficit bajó de 47.4\,\% a 42.2\,\%, porque el ingreso
laboral creció al doble de velocidad que el consumo.
El origen de ese déficit agregado importa: no es un resultado de la imputación, sino de la
contabilidad nacional. Los dos términos de la identidad tienen agregado de control y por tanto
su diferencia queda determinada antes de repartir nada por edad. El ingreso laboral es la
remuneración de asalariados más dos tercios del ingreso mixto neto de depreciación,
ecuación~\eqref{eq:ylmacro}; el consumo es el gasto final de hogares y gobierno neto de impuestos
a los productos, ecuación~\eqref{eq:cmacro}. En México el primero representa entre 39.1\,\% y
44.4\,\% del producto interno bruto y el segundo entre 72.9\,\% y 78.4\,\%, de modo que la
diferencia equivale a entre 31.5\,\% y 35.8\,\% del producto. Esa brecha existe en las cuentas
nacionales con independencia de cualquier supuesto de este documento; el ejercicio sólo decide
cómo se distribuye por edad, no cuánto vale.
\begin{equation}
\label{eq:cmacro}
C(t) = C^{hog}(t) + C^{gob}(t) - T^{prod}(t),
\qquad
LCD(t) = C(t) - Y_L(t)
\end{equation}
La participación laboral relativamente baja del producto en México fija así un déficit agregado
alto que ninguna variante de las imputaciones puede cerrar. Lo que sí depende del procedimiento
es en qué edades se concentra, y esa es la pregunta que aborda el resto de la sección.
El perfil por edad del consumo es mucho más plano que el del ingreso laboral, y su distancia entre
extremos se amplió. En términos reales el consumo per cápita de los menores de 15 años cayó
20.3\,\%, de 30\,364 a 24\,209 pesos trimestrales, mientras el de los mayores de 60 subió
7.3\,\%, de 45\,344 a 48\,666. La razón entre ambos extremos pasó de 1.49 a 2.01, y la del grupo
de 25 a 59 años frente a la niñez de 1.86 a 2.54.
El déficit per cápita bajó en la niñez y en la juventud y subió en la vejez. Entre 2016 y 2024 el
déficit real cayó 20.5\,\% entre los menores de 15 años y 17.9\,\% entre los jóvenes de 15 a 24,
mientras subió 71.7\,\% a partir de los 60 años, de 18\,569 a 31\,879 pesos trimestrales; en el
grupo de 25 a 59 bajó 5.9\,\%, de 11\,601 a 10\,914. La ventana de edades con superávit se cerró
en paralelo: abarcaba diez años en 2020 --- de los 40 a los 49 --- y en 2022 y 2024 no hay ninguna
edad en la que el ingreso laboral per cápita del conjunto de la población supere a su consumo.
\begin{figure}[H]
\centering
\includegraphics[width=\linewidth]{fig_c_lcd.pdf}
\caption{Consumo, ingreso laboral y déficit del ciclo de vida por edad. Panel izquierdo: consumo
e ingreso laboral en 2024, con el área sombreada según el signo de la diferencia. Panel derecho:
déficit por el primer método, 2016--2024. Pesos constantes de 2024, trimestrales per cápita.}
\label{fig:c_lcd}
\end{figure}
Ese déficit que no se cierra en ninguna edad esconde dos ciclos de vida distintos. Separado por
sexo, el perfil masculino sí completa el ciclo: en 2024 el ingreso laboral de los hombres supera a
su consumo entre los 33 y los 62 años, una ventana de treinta edades cuyo punto más alto son
18\,547 pesos trimestrales per cápita de superávit hacia los 47 años. El perfil femenino no cruza
el eje en ninguna edad de ningún año de la serie: su valor más bajo en edad activa, 18\,700 pesos
hacia los 45 años, sigue siendo un déficit. El agregado deficitario en todo el ciclo es, por
tanto, resultado de composición, no una característica del ciclo de vida mexicano.
La distancia entre ambos perfiles se concentra en la edad activa y no se cerró durante el periodo.
Entre los 25 y los 59 años los hombres registran en 2024 un superávit per cápita de 7\,033 pesos
trimestrales y las mujeres un déficit de 26\,461, y la misma configuración se repite en los cinco
levantamientos. Las mujeres concentran 76.3\,\% del déficit agregado de 2024, frente a 74.5\,\%
en 2016, de modo que financiar el ciclo de vida en México es, en su mayor parte, financiar el
déficit femenino.
\begin{figure}[H]
\centering
\includegraphics[width=\linewidth]{fig_lcd_sexo.pdf}
\caption{Déficit del ciclo de vida por sexo. Panel izquierdo: perfil por edad en 2024, con el área
sombreada sobre las edades en que el ingreso laboral masculino supera al consumo. Panel derecho:
saldo per cápita de la población de 25 a 59 años, 2016--2024; los valores negativos son superávit.
Pesos constantes de 2024, trimestrales per cápita.}
\label{fig:lcd_sexo}
\end{figure}
\subsection{Financiamiento por rangos de edad}
La composición del financiamiento se invierte entre la niñez y la vejez. En 2024 los componentes
que financian el déficit de los menores de 15 años se reparten en 57.4\,\% de reasignación por
activos y 31.2\,\% de transferencias privadas entre hogares, con 8.8\,\% de transferencias
públicas netas. En el grupo de 60 años o más la jerarquía se invierte: las transferencias públicas
netas aportan 54.5\,\% y la reasignación por activos 42.7\,\%. El contraste sitúa a la niñez como
un déficit financiado por vía privada y a la vejez como uno financiado por vía pública.
El grupo de 25 a 59 años es contribuyente neto tanto al sistema público como al circuito privado.
Su pago neto de transferencias públicas asciende a 12\,960 pesos trimestrales per cápita en 2024
y su saldo de transferencias privadas entre hogares es negativo en 2\,145 pesos. Ese pago público
neto creció 39.1\,\% en términos reales respecto de 2016, cuando era de 9\,317 pesos trimestrales
per cápita. El aumento no proviene de una mayor carga tributaria sobre el grupo, sino de lo que
dejó de recibir: su pago bruto per cápita subió 15\,\% en términos reales, mientras las
transferencias públicas que percibe cayeron a la mitad, porque su participación en el total
transferido pasó de 27.7\,\% a 13.4\,\%.
La transferencia pública neta por adulto mayor se mantuvo estable en términos reales mientras la
contribución que la financia creció con rapidez. Pasó de 36\,543 pesos trimestrales per cápita en
2016 a 37\,943 en 2024, un aumento real de 3.8\,\%, frente al 39.1\,\% que subió el pago neto
por persona en edad activa. El sistema desplazó transferencias hacia las edades mayores.
\begin{figure}[H]
\centering
\includegraphics[width=\linewidth]{fig_financiamiento_edad.pdf}
\caption{Financiamiento del déficit por grupo de edad. Panel izquierdo: componentes en 2024.
Panel derecho: déficit per cápita por grupo de edad, 2016--2024. Pesos constantes de 2024,
trimestrales per cápita.}
\label{fig:fin_edad}
\end{figure}
\subsection{Financiamiento de jóvenes, mujeres y adultos mayores}
Entre los jóvenes de 15 a 24 años el déficit de las mujeres supera al de los hombres en 36\,\%.
En 2024 el déficit per cápita es de 39\,551 pesos trimestrales para ellas y de 29\,133 para
ellos. La diferencia se financia sobre todo por vía privada: las mujeres jóvenes reciben 9\,045
pesos en transferencias entre hogares frente a 6\,232 de los hombres, y 1\,958 pesos en
transferencias del exterior frente a 329, casi seis veces más.
Entre los adultos mayores la diferencia por sexo es mayor y su composición es distinta. El
déficit per cápita de las mujeres de 60 años o más es de 42\,322 pesos trimestrales, 2.2 veces el
de los hombres del mismo grupo, que es de 19\,109. Dos de los tres canales operan en su contra:
las mujeres reciben 33\,224 pesos de transferencias públicas netas frente a 43\,713 de los
hombres, una diferencia de 24.0\,\%, y 27\,414 por reasignación de activos frente a 32\,648. El
único canal que las favorece es el privado, con un saldo positivo de 748 pesos frente a uno
negativo de 9\,082 en el caso de los hombres, insuficiente para compensar los otros dos.
El perfil femenino mantiene déficit positivo en los cuatro grupos de edad. En el grupo de 25 a 59
años, donde el conjunto de la población registra un déficit per cápita de 10\,914 pesos
trimestrales, el de las mujeres asciende a 26\,461. Las transferencias del exterior que reciben
las mujeres de esa edad son de 4\,185 pesos trimestrales, 1.7 veces el promedio del grupo, lo que
sitúa a las remesas como un canal de financiamiento con sesgo de sexo marcado en la edad activa.
El nivel de ese canal depende del reescalamiento (véase el Cuadro~\ref{tab:remesas}); la razón
entre sexos, en cambio, es la misma con y sin ajuste, porque el factor es único por año y no
altera la forma relativa del perfil.
\begin{figure}[H]
\centering
\includegraphics[width=\linewidth]{fig_financiamiento_grupos.pdf}
\caption{Componentes del financiamiento en 2024 para tres subpoblaciones: jóvenes de 15 a 24 años
por sexo, adultos mayores de 60 años o más por sexo, y mujeres por grupo de edad. Pesos constantes de 2024,
trimestrales per cápita.}
\label{fig:fin_grupos}
\end{figure}
\subsection{Validación}
\label{sec:validacion}
El procedimiento se sometió a tres comprobaciones reproducibles. La primera verifica la
neutralidad de las imputaciones respecto de los agregados de control: para las siete variables
del mapa de escalamiento y los cinco años, el agregado micro ajustado reproduce el agregado de
control dentro de la tolerancia declarada, sin diferencias relevantes. La segunda comprueba que
la agregación por edad no pierda información respecto de la base persona-año, con el mismo
resultado. La tercera contrasta la corrida contra una salida de referencia y confirma que
ninguna variable con agregado de control asociado modificó su total anual.
La identidad contable se contrastó después por los dos métodos. El primero cierra contra el
agregado de control, porque sus dos componentes tienen correspondencia directa. El segundo no
cierra, y la brecha no es un residuo desconocido: todos sus términos excepto el bloque privado
tienen agregado de control y agregan a él, de modo que la diferencia equivale al agregado de las
transferencias observadas entre hogares. En 2024 esa brecha es de 277 mil millones de pesos
anuales, frente a un agregado de transferencias entre hogares de 277 mil millones, con un residuo
de 61 millones, equivalente a 0.001\,\% del déficit total. La brecha se reporta y no se fuerza a
cero: un cierre forzado exigiría un agregado de control inexistente, o bien el reescalamiento del
bloque privado contra un objetivo nulo, que lo eliminaría.
Un tercer contraste, de naturaleza distinta, aplica al bloque de transferencias del exterior. Su
factor de reescalamiento es el mayor del ejercicio y la base de receptores observados es estrecha,
por lo que los resultados que dependen de él se reportan con y sin el ajuste. El
Cuadro~\ref{tab:remesas} lo hace. El canal pasa de financiar 2.6\,\% del déficit del grupo de 25
a 59 años a financiar 23.2\,\%, y de 1.9\,\% a 17.6\,\% en el grupo de 60 años o más; en la
niñez el peso es marginal bajo cualquiera de las dos lecturas. La composición interna del canal no
se ve afectada: al ser un factor único por año, la razón entre sexos en la edad activa es 6.8 en
ambos casos.
\input{tables/tab_remesas.tex}
La lectura correcta es que la \emph{forma} del canal por edad y por sexo proviene de la encuesta y
es robusta al ajuste, mientras que su \emph{nivel} proviene del agregado de control y descansa en
el supuesto de que toda la subcaptación ocurre en el monto declarado y ninguna en el número de
hogares receptores. Corregir esa asimetría exige repartir el ajuste entre los dos márgenes: imputar
receptores adicionales con un modelo de propensión --- sobre entidad federativa, composición del
hogar y posición en la distribución del ingreso --- y escalar los montos por un factor
sustancialmente menor, con un objetivo externo de incidencia tomado de las estadísticas de envíos.
Ese refinamiento queda fuera del alcance de este documento y se anota como extensión pendiente.
\subsection{Implicaciones}
El resultado con mayor consecuencia para el espacio fiscal es la asimetría en el financiamiento
de los dos extremos del ciclo. La niñez se financia en 88.6\,\% con recursos privados --- activos
y transferencias entre hogares --- mientras la vejez se financia en 54.5\,\% con transferencias
públicas netas. Un cambio demográfico que aumente el peso relativo de los mayores de 60 años se
traduce así de forma directa en presión sobre el presupuesto público, mientras que un cambio en
el peso de la niñez se absorbe sobre todo dentro de los hogares.
La segunda implicación proviene de la divergencia entre lo que recibe cada adulto mayor y lo que
aporta cada persona en edad activa. El déficit per cápita de los mayores de 60 años creció
71.7\,\% en términos reales entre 2016 y 2024, mientras la transferencia pública neta que reciben
subió 3.8\,\% y la contribución neta por persona en edad activa subió 39.1\,\%. El peso relativo
del grupo de 60 años o más pasó de 11.3\,\% a 15.1\,\% de la población en el mismo lapso, de modo
que la carga por aportante creció sin que la prestación real por beneficiario avanzara, y el resto
del déficit se trasladó a la reasignación por activos y a los hogares.
La tercera se refiere a la desigualdad horizontal por sexo. El déficit de las mujeres de 60 años
o más duplica al de los hombres, y las transferencias públicas netas que reciben son 24.0\,\%
menores, de modo que el sistema público no compensa la brecha sino que se suma a ella. El mismo
patrón aparece en la edad activa, como se mostró antes: el déficit femenino más que duplica al
del grupo completo. La incidencia del gasto público por sexo, evaluada sobre estos perfiles, no
corrige un punto de partida desigual que la política pública no genera pero tampoco revierte.
% ---------------------------------------------------------------------------
\section{Conclusiones}
\label{sec:conclusiones}
% ---------------------------------------------------------------------------

La contribución central de este documento es metodológica: una base de Cuentas
Nacionales de Transferencia para México construida con cierre verificable contra el
Sistema de Cuentas Nacionales, documentación explícita de cada imputación y sus
supuestos, y reporte de la brecha entre los dos métodos de cálculo del déficit del
ciclo de vida en lugar de un cierre forzado. Esa transparencia es lo que hace
comparables los resultados con otros países y reproducible el ejercicio ante cambios
en los datos o en las decisiones metodológicas.

Sobre esa base emergen tres hallazgos sustantivos. El primero es que el déficit
agregado del ciclo de vida en México no crece, pero su composición cambia: bajo la
estructura actual de transferencias y perfiles por edad, el déficit se desplaza hacia
la vejez al mismo ritmo en que envejece la población. Esto importa fiscalmente porque
los dos extremos del ciclo se financian por vías opuestas: la niñez carga el costo
sobre el circuito privado --- activos y transferencias entre hogares ---, mientras la
vejez depende en más de la mitad de transferencias públicas netas. Manteniendo
constantes los perfiles observados, cada punto porcentual que gane el grupo de mayores
en la distribución poblacional se traduce en presión presupuestaria directa; un cambio
equivalente en el peso de la niñez se absorbe dentro de los hogares.

El segundo hallazgo es que el promedio nacional oculta dos ciclos de vida
estructuralmente distintos por sexo. En el marco monetario de las NTA, el perfil
masculino alcanza superávit en un tramo amplio de la edad activa; el femenino no lo
alcanza en ninguna edad de ninguno de los cinco levantamientos. Esta asimetría debe
leerse dentro de los límites del instrumento: las NTA monetarias no incorporan el
trabajo no remunerado, de modo que el déficit femenino observado no es una medida
integral de la contribución económica de las mujeres, sino de su posición dentro del
ciclo de vida monetario tal como lo registran la encuesta y las cuentas nacionales.
Con esa cautela explícita, el resultado es informativo: el sistema público no compensa
la brecha de financiamiento por sexo en la vejez, y las transferencias privadas entre
hogares --- el único canal que favorece a las mujeres mayores frente a los hombres ---
resultan insuficientes para cerrarla.

El tercer hallazgo concierne a la dinámica interna del sistema público entre 2016 y
2024. La carga por aportante en edad activa creció de forma sustantiva, no por un
aumento proporcional de su contribución bruta, sino porque su participación en las
transferencias públicas recibidas cayó a la mitad. El sistema desplazó recursos hacia
las edades mayores sin ampliar la base que los financia, mientras la prestación real
por beneficiario mayor apenas varió. Esta configuración es descriptiva de la estructura
observada, no una proyección: cuánto persiste bajo distintos escenarios de política,
formalización o comportamiento laboral es una pregunta que requiere un modelo de
segunda etapa.

La base construida aquí es ese punto de partida. Los perfiles por edad, sexo y posición
en la distribución del ingreso que produce están listos para alimentar simulaciones
demográficas y fiscales que distingan los efectos de composición poblacional de los
efectos de política, y para extenderse hacia las Cuentas Nacionales de Inclusión
mediante la desagregación socioeconómica. Esas extensiones constituyen la segunda
etapa analítica.

% ---------------------------------------------------------------------------
\appendix
\section{Series del Sistema de Cuentas Nacionales}
\label{anexo:series}
% ---------------------------------------------------------------------------
\input{tables/tab_series_macro.tex}
\FloatBarrier
% ---------------------------------------------------------------------------
\section{Correspondencia entre variables micro y agregados de control}
\label{anexo:mapeo}
% ---------------------------------------------------------------------------
\input{tables/tab_scale_map.tex}
\input{tables/tab_excluidas.tex}
\FloatBarrier
% ---------------------------------------------------------------------------
\section{Registro de supuestos}
\label{anexo:supuestos}
% ---------------------------------------------------------------------------
\input{tables/tab_supuestos.tex}
\FloatBarrier
% ---------------------------------------------------------------------------
\section{Factores de reescalamiento}
\label{anexo:lambda}
% ---------------------------------------------------------------------------
\input{tables/tab_lambda.tex}
\FloatBarrier
% ---------------------------------------------------------------------------
\section{Limitaciones y puntos a verificar}
\label{anexo:limitaciones}
% ---------------------------------------------------------------------------
\input{tables/tab_limitaciones.tex}
\FloatBarrier
\noindent El cuadro recoge sólo los puntos abiertos. Dos que aparecían en versiones previas de
este registro ya están resueltos y sus correcciones están incorporadas a las cifras: la escala de
equivalencia por edad, que se estimaba sobre la muestra sin expandir, y el reescalamiento del
consumo, que escalaba una segunda vez el bloque de gobierno ya calibrado. Ninguna de las dos
alteraba los agregados anuales; ambas afectaban la forma por edad de los perfiles.
% ---------------------------------------------------------------------------
\bibliographystyle{plainnat}
\bibliography{referencias}
\end{document}